% Kompilacja: pdflatex main && bibtex main && pdflatex main && pdflatex main
% Dane liczbowe (definicje.tex), tabele (tabele/*.csv) i rysunki (rysunki/) generuje
% analiza_odpowiedzi_naukowa.py; tekst (main.tex, sekcja-*.tex) jest edytowany ręcznie.
\documentclass[11pt]{article}
\usepackage[utf8]{inputenc}
\usepackage[T1]{fontenc}
\usepackage[polish]{babel}
\usepackage[margin=2.5cm]{geometry}
\usepackage{booktabs,graphicx,xcolor,amsmath}
\usepackage[hidelinks]{hyperref}
\usepackage{float}
\usepackage{changepage} % do marginesow
\usepackage{enumitem}% numeracja rzymska listy definicji
\usepackage{doi}
% BIBLIOGRAFIA:
\usepackage[numbers,sort&compress]{natbib}
\bibliographystyle{plainnat}

\extrafloats{100}
\newcommand{\fillin}[1]{\textcolor{red}{[#1]}}
\input{definicje.tex}

\usepackage{authblk}
\renewcommand\Authand{ i }\renewcommand\Authands{ i }
\author[1,*]{Bartosz Lewandowski \small{(ORCID:0009-0004-1158-0693)}}
\author[2]{Monika Szerszeń \small{(ORCID: 0009-0009-6625-3979)}}
\author[3]{Magdalena Dowgalec \small{(ORCID:0009-0005-4469-2832)}}
\author[4]{Ewaryst Tkacz \small{(ORCID:0000-0003-2580-7954)}}
\author[5]{Maciej J. Wróbel \small{(ORCID:0000-0003-1865-2122)}}

\affil[1,4]{Department of Clinical Engineering, Academy of Silesia, Katowice, Poland}
\affil[2]{Banking Process Support and Optimization Department, Powiślański Cooperative Bank, Kwidzyn, Poland}
\affil[3]{Legal Department, Atende SA, Warsaw, Poland}% SPRAWDŹ: afiliacja M. Szerszeń (wcześniej podwójny numer [4])
\affil[5]{Department of Otolaryngology and Oncological Laryngology, Collegium Medicum NCU, Toruń, Poland}

% Autor korespondencyjny
\affil[*]{Autor korespondencyjny: dr inż. Bartosz Lewandowski, bartosz.lewandowski@akademiaslaska.pl}

\title{Poprawność odpowiedzi pracowników w powtarzanych pytaniach: zapominanie i utrwalone błędy w trzyletnim programie mikrouczenia}


\date{}
\begin{document}
\maketitle


\begin{abstract}
\noindent\textbf{Wprowadzenie.} Powtórki rozłożone w czasie poprawiają zapamiętywanie, jednak
skuteczność mikrouczenia w organizacjach rzadko oceniano na podstawie wieloletnich danych
o odpowiedziach pracowników.
\textbf{Metody.} Przeanalizowano \NResponses{} odpowiedzi \NEmployees{} pracowników banku na
\NItems{} pytania wielokrotnego wyboru, udzielonych od \DateStart{} do \DateEnd. Pytania wracały
według harmonogramu niezależnego od wyniku odpowiedzi (mediana odstępu \GapMedian{} dni), a po każdej odpowiedzi wyświetlano poprawną odpowiedź
z wyjaśnieniem. Poprawność analizowano w kolejnych próbach tego samego pytania oraz według wyniku
poprzedniej próby i odstępu.
\textbf{Wyniki.} Poprawność w pierwszej próbie wynosiła \FirstAcc\%. W sekwencjach liczących co
najmniej \CohortK{} prób wzrosła z \CohortAccOne\% do \CohortAccK\% (o \SeqChange{} p.p., 95\% PU
\SeqChangeCI). Korekt błędów było więcej niż przypadków zapominania (\NCorrections{} wobec
\NForgetting).
Odsetek zapominania rósł wraz z odstępem i malał wraz z liczbą poprawnych odpowiedzi z rzędu,
z \ForgetRunOne\% po jednej do \ForgetRunMax\% po co najmniej \ForgetRunK{}. Powtórzone błędy
w \SameDistractorPct\% przypadków dotyczyły tego samego dystraktora (oczekiwane przy wyborze zgodnym
z popularnością dystraktorów \SameDistractorExpectedPop\%), niezależnie od odstępu.
\textbf{Wnioski.} Poprawność odpowiedzi na powracające pytania rosła, lecz wiedza wymagała dalszych
powtórek, a część błędów utrzymywała się mimo informacji zwrotnej. Harmonogram dopasowany do wyników
pracownika i informacja zwrotna odnosząca się do wybranego dystraktora mogą zwiększyć skuteczność
takich programów.

\medskip\noindent\textbf{Słowa kluczowe:} powtórki rozłożone w czasie; mikrouczenie; szkolenia
pracownicze; informacja zwrotna; zapominanie; analityka uczenia się
\end{abstract}

% =============================================================================
% Wprowadzenie. Plik edytowany ręcznie -- skrypt analiza_odpowiedzi_naukowa.py go NIE nadpisuje.
% Liczby: definicje.tex; tabele: \TabDir (CSV); rysunki: \FigDir; literatura: references.bib.
% =============================================================================
\section{Wprowadzenie}

Przypominanie informacji z pamięci utrwala je silniej niż ponowne czytanie (efekt testowania),
a nauka rozłożona w czasie jest trwalsza niż nauka skumulowana (efekt rozłożenia w czasie)
\citep{roediger2006,cepeda2006}. Oba efekty łączy edukacja rozłożona w czasie (ang. \emph{spaced
education}), w której krótkie pytania wielokrotnego wyboru wracają cyklicznie przez tygodnie lub
miesiące. W badaniach z udziałem lekarzy i innych pracowników ochrony zdrowia poprawiała ona
przyswajanie i utrwalanie wiedzy \citep{kerfoot2009,phillips2019}.

W organizacjach podobną formę ma \emph{mikrouczenie} (ang. \emph{microlearning}), czyli krótkie
jednostki szkoleniowe, często w formie quizów. Dowody jego skuteczności pochodzą jednak głównie
z badań na małych próbach i z tzw. szarej literatury \citep{taylor2022}, a efekty obowiązkowych
szkoleń pracowniczych, np.\ z bezpieczeństwa informacji, bywają niewielkie lub zanikają z czasem
\citep{reinheimer2020}. Skuteczność szkoleń ocenia się zwykle przez transfer, czyli zastosowanie
wiedzy w pracy \citep{blume2010}. Transfer wymaga jednak, aby wiedza została najpierw nabyta
i utrwalona.

Ten wcześniejszy etap można sprawdzić na danych z platform mikrouczenia. Platforma zapisuje każdą
odpowiedź wraz z czasem i wynikiem, a to samo pytanie wraca do tego samego pracownika wielokrotnie.
Kolejne odpowiedzi stanowią więc powtarzane pomiary tej samej wiedzy u tej samej osoby. Podobne dane
wykorzystuje się w modelach ucznia, stosowanych w systemach tutorskich i aplikacjach do nauki języków
\citep{corbett1994,settles2016}.

Badanie przeprowadzono w banku, w którym pracownicy odpowiadali na krótkie pytania wyświetlane na
komputerze w trakcie dnia pracy, a system gromadził odpowiedzi przez trzy lata. Pytania wracały według
harmonogramu niezależnego od wyniku odpowiedzi, a po każdej odpowiedzi wyświetlano odpowiedź poprawną.
Taki harmonogram pozwala porównać korekty błędów i zapominanie bez selekcji, którą wprowadzają
algorytmy adaptacyjne, co w danych z platform szkoleniowych występuje rzadko. Badania laboratoryjne
wskazują ponadto, że wybrany dystraktor może utrwalić się w pamięci, a błędy popełnione z wysoką
pewnością mogą powracać mimo informacji zwrotnej \citep{butler2008,butler2011}. Nie wiadomo, czy
zjawisko to występuje w szkoleniach pracowniczych. Sformułowano trzy pytania badawcze:
(PB1) czy odsetek poprawnych odpowiedzi rośnie, gdy to samo pytanie wraca do tego samego pracownika;
(PB2) jak poprawność zależy od wyniku poprzedniej próby i od czasu, który od niej upłynął;
(PB3) czy powtórzone błędy utrzymują się mimo informacji zwrotnej, co wskazywałoby raczej na błędne
przekonania niż na zgadywanie.

% =============================================================================
% Materiał i metody. Plik edytowany ręcznie -- skrypt analiza_odpowiedzi_naukowa.py go NIE nadpisuje.
% Liczby: definicje.tex; tabele: \TabDir (CSV); rysunki: \FigDir; literatura: references.bib.
% Analizy uzupełniające (zmiana na poziomie pracowników, analiza pytań) opisano w suplement.tex.
% =============================================================================
\section{Materiał i metody}

\subsection{Program szkoleniowy}
Badanie przeprowadzono w banku. Pytania przygotował \textit{Wydział wsparcia i optymalizacji procesów
bankowych}, który decydował również o udziale pracowników w programie. Pytania dotyczyły wiedzy
wykorzystywanej w codziennej pracy w banku. Każde było pytaniem wielokrotnego wyboru z trzema do pięciu
opcjami odpowiedzi, z których jedna była poprawna, a pozostałe stanowiły \emph{dystraktory}.
Pytania, opcje odpowiedzi i wyjaśnienia sformułowano w języku polskim. 

Aplikacja szkoleniowa uruchamiała się automatycznie w dni robocze na komputerach pracowników
i wyświetlała od jednego do czterech losowo wybranych pytań dziennie, w godzinach 9:00--16:00.
Po każdej odpowiedzi, poprawnej lub błędnej, pokazywała poprawną odpowiedź z wyjaśnieniem (\emph{informacja zwrotna}). To samo pytanie
mogło wrócić do pracownika najwcześniej po siedmiu dniach, a harmonogram powtórek nie zależał od wyniku
odpowiedzi. Na czas przeglądu i aktualizacji pytań program zawieszano. W analizowanych danych mediana
liczby odpowiedzi pracownika w dniu aktywności wynosiła \DailyMedian{} (IQR \DailyQOne--\DailyQThree).
To samo pytanie wracało po medianie \GapMedian{} dni (IQR \GapQOne--\GapQThree).

Przykładowe pytanie zadanie uczestnikowi programu, z zakresu cyberbezpieczeństwa brzmiało:
,,Która z poniższych sytuacji jest phishingiem?''. Opcje odpowiedzi: 
,,1 - Atak typu DDoS'',
,,2 - Instalacja oprogramowania antywirusowego'',
,,3 - Tworzenie kopii zapasowych danych'',
,,4 - Fałszywy e-mail z linkiem do strony logowania banku (oznaczona jako odpowiedź poprawna)''.
Wyjaśnienie wyświetlane po odpowiedzi:
,,Phishing to metoda oszustwa, w której ktoś podszywa się pod zaufaną osobę lub instytucję (np. bank, kuriera, urząd), żeby wyłudzić od Ciebie poufne dane, takie jak hasła, numery kart czy loginy. Najczęściej wygląda to jak e-mail z "banku" z pilną prośbą o zalogowanie, SMS z linkiem do "dopłaty za paczkę" albo fałszywa strona łudząco podobna do prawdziwej. Kluczowy mechanizm polega na tym, że oszust wywołuje pośpiech lub stres, żebyś kliknął/a bez zastanowienia. WAŻNE: nie klikaj w podejrzane linki, zawsze sprawdzaj adres strony (URL), nie podawaj danych logowania przez link z wiadomości, a jeśli coś wygląda na "zbyt pilne", to często oszustwo.''.

\subsection{Dane i grupa badana}
\label{sec:dane}
Analizowano zarejestrowane odpowiedzi pracowników. Każdy rekord zawierał pytanie, wybraną opcję,
kod pracownika, datę i godzinę odpowiedzi oraz informację, czy odpowiedź była poprawna. Zbiór analityczny liczył \NResponses{} odpowiedzi
% Po wykluczeniu odpowiedzi testowych, pytań z jedną zaobserwowaną opcją oraz powtórek niezgodnych
% z regułą siedmiu dni zbiór analityczny liczył \NResponses{} odpowiedzi
(tab.~\ref{tab:zbior_danych}). Bank przekazał ponadto zestawienie wieku i płci pracowników.

\begin{table}[H] 
\centering
\caption{Charakterystyka zbioru danych.}
\label{tab:zbior_danych} 
\small
\tabelacsv{lr}{zbior_danych}
\uwagatabeli{IQR: rozstęp międzykwartylowy. Sekwencja to seria prób jednego pracownika dotyczących
jednego pytania. Liczba zaobserwowanych opcji odpowiedzi bywa mniejsza niż liczba opcji w pytaniu,
ponieważ części opcji nie wybrał żaden pracownik. Wykluczone odpowiedzi testowe: \NExcluded. Wykluczone
pytania z jedną zaobserwowaną opcją odpowiedzi: \NExcludedItems{} (\NExcludedItemResp{} odpowiedzi).
Wykluczone powtórki udzielone wcześniej niż po \MinGapDays{} dniach: \NExcludedShortRep.}
\end{table}

Dane demograficzne były dostępne dla \DemoMatched{} ze \DemoN{} pracowników (\DemoMatchedPct\%).
Kobiety stanowiły \DemoWomenPct\% osób o znanej płci ($n = \DemoWomen$). Mediana wieku wynosiła
\AgeMedian{} lat (IQR \AgeQOne--\AgeQThree; zakres \AgeMin--\AgeMax), a najliczniejszy był przedział
\AgeModalGroup{} lat (\AgeModalGroupPct\%; tab.~\ref{tab:demografia}). Wiek uwzględniono w analizie
pierwszej próby i tempa wzrostu poprawności; płci nie analizowano.
% SPRAWDŹ: przedział modalny i proporcje płci po zmianie danych

\begin{table}[H]
\centering
\caption{Charakterystyka demograficzna pracowników objętych analizą.}
\label{tab:demografia}
\small
\tabelacsv{lr}{demografia}
\uwagatabeli{Odsetki płci liczono wśród osób o znanej płci, a odsetki przedziałów wieku wśród osób
o znanym wieku ($n = \DemoAgeN$). Brak danych dotyczy głównie pracowników nieobecnych w zestawieniu
demograficznym (\DemoUnmatched). IQR: rozstęp międzykwartylowy; SD: odchylenie standardowe.}
\end{table} 

W bazie danych odpowiedzi, pracowników oznaczono losowym kodem, a zestawienie demograficzne połączono z zarejestrowanymi odpowiedziami według
tego kodu. Tego rodzaju kodu nie można odtworzyć przy znajomości listy pracowników, dlatego analizowane dane
stanowiły dane w pełni zanonimizowane.

\subsection{Analiza statystyczna}
\label{sec:analiza}
Odpowiedzi jednego pracownika na jedno pytanie, uporządkowane w czasie, tworzyły \emph{sekwencję}.
Pozycję odpowiedzi w sekwencji nazwano \emph{numerem próby}, a pierwszą odpowiedź, czyli pierwszy
kontakt pracownika z pytaniem, \emph{pierwszą próbą}. \emph{Odstęp} oznaczał liczbę dni od poprzedniej
próby w tej samej sekwencji, a \emph{poprawność} odsetek poprawnych odpowiedzi. Poprawność podawano
z 95-procentowymi przedziałami ufności (PU) wyznaczonymi metodą bootstrap: pracowników losowano ze
zwracaniem (\BootReps{} replikacji), co uwzględnia zależność odpowiedzi udzielanych przez tę samą osobę.
W modelach regresji błędy standardowe klastrowano względem pracowników.

\emph{Pierwsza próba.} Poprawność w pierwszej próbie porównano między przedziałami wieku. Zależność od
wieku (podział grup po 10 lat) oszacowano regresją logistyczną z efektami stałymi pytania, czyli z osobnym wyrazem
wolnym dla każdego pytania. Trudność pytań nie była oznaczona w danych, dlatego efekty stałe szacowano
z odpowiedzi wszystkich pracowników. Model porównywał więc pracowników w różnym wieku odpowiadających
na te same pytania. Dla porównania dopasowano go również bez efektów stałych pytania.

\emph{Zmiany w sekwencji (PB1).} Poprawność obliczono dla kolejnych numerów prób; próby
\MaxAttemptPooled{} i dalsze połączono. Z każdą próbą ubywało sekwencji, a pracownicy kontynuujący
odpowiedzi mogli różnić się od pozostałych, co zniekształca krzywe zagregowane \citep{murray2013}.
Zmianę oceniono więc także w \emph{kohorcie}, czyli w sekwencjach liczących co najmniej \CohortK{}
prób, ograniczonych do prób od 1 do \CohortK{}. Poprawność w każdej z tych prób obliczano dla tych
samych sekwencji, a miarą zmiany była różnica między próbą~\CohortK{} a próbą~1 w punktach
procentowych (p.p.). 
Związek tempa wzrostu z wiekiem sprawdzono w kohorcie regresją logistyczną z interakcją $\log_2$
numeru próby i wieku oraz z efektami stałymi pracownika i pytania.

\emph{Wynik poprzedniej próby i odstęp (PB2).} Każdą próbę od drugiej wzwyż przypisano do jednej
z czterech kategorii według wyniku tej i poprzedniej próby: odpowiedź poprawna po poprawnej,
\emph{korekta błędu} (poprawna po błędnej), \emph{zapominanie} (błędna po poprawnej) albo
\emph{powtórzony błąd} (błędna po błędnej). Ponieważ odsetek korekt błędów ($r$) i odsetek
zapominania ($f$) mają różne mianowniki, porównano liczby obu zdarzeń oraz wyznaczono poziom równowagi
$r/(r+f)$, do którego dążyłaby poprawność przy stałych odsetkach. Odsetki porównano w sześciu
przedziałach odstępu: \MinGapDays--13, 14--29, 30--59, 60--119, 120--239 oraz co najmniej 240 dni. Odsetek zapominania zestawiono także według liczby
poprawnych odpowiedzi z rzędu przed daną próbą oraz osobno po korekcie błędu. Odpowiedź błędną
w próbie następującej po korekcie błędu określono jako \emph{nawrót błędu} i sprawdzono, czy pracownik
wybrał w niej ten sam dystraktor co przed korektą. Sprawdzono również, czy harmonogram
powtórek zależał od wyniku poprzedniej próby.

\emph{Utrwalone błędy (PB3).} Poprawność porównano według liczby \emph{błędów z rzędu}, czyli
kolejnych odpowiedzi błędnych bezpośrednio poprzedzających daną próbę, od 0 do co najmniej \RunMax.
Dla dwóch kolejnych błędów obliczono, jak często pracownik wybierał ten sam dystraktor, i porównano
ten odsetek z oczekiwanym przy wyborze losowym oraz przy wyborze zgodnym z popularnością dystraktorów
w danym pytaniu. Odsetek ten zestawiono także w trzech przedziałach odstępu: krótszym niż 30 dni,
30--89 dni oraz co najmniej 90 dni. Obliczono też, jaką część powtórzonych błędów skupia 10\% pytań, w których było ich
najwięcej.

Obliczenia wykonano w języku Python \VerPython{} z użyciem bibliotek pandas \VerPandas, NumPy
\VerNumpy, SciPy \VerScipy, statsmodels \VerStatsmodels{} i Matplotlib \VerMatplotlib.

% \subsection{Analiza opisowa}
% Odsetki poprawnych odpowiedzi podano z 95-procentowymi przedziałami ufności (PU) z bootstrapu
% klastrowego, w którym losowano pracowników ze zwracaniem (\BootReps{} replikacji); procedura ta
% uwzględnia zależność odpowiedzi udzielanych przez tę samą osobę. Ponieważ osoby kontynuujące do
% kolejnych prób mogą różnić się od tych, które przestają, zagregowane krzywe uczenia mogą być
% zniekształcone przez ubywanie obserwacji \citep{murray2013,goutte2018}; krzywą obliczono więc
% dodatkowo dla zbalansowanej kohorty sekwencji liczących co najmniej \CohortK{} prób. Zależność odsetka
% błędów od numeru próby opisano funkcją potęgową \citep{newell1981}. Ponieważ harmonogram powtórek nie
% był adaptacyjny, sprawdzono, czy ponowne zadanie pytania zależało od wyniku poprzedniej próby.

% \subsection{Modele statystyczne}
% Niech $Y_{ijt}$ oznacza poprawność ($1$ albo $0$) odpowiedzi pracownika $i$ na pytanie $j$ w próbie
% $t$. Zastosowano trzy modele regresji logistycznej o wspólnej postaci
% \begin{equation*}
% \operatorname{logit} P(Y_{ijt}=1) = \alpha_i + \gamma_j + \mathbf{x}_{ijt}^{\top}\boldsymbol{\beta},
% \end{equation*}
% gdzie $\alpha_i$ i $\gamma_j$ są efektami stałymi pracownika i pytania. Efekty stałe pochłaniają
% stałe różnice poziomu wiedzy między osobami i trudności między pytaniami, dlatego współczynniki
% $\boldsymbol{\beta}$ opisują zmiany w obrębie tej samej osoby i tego samego pytania. Modele różniły
% się zbiorem odpowiedzi i wektorem predyktorów $\mathbf{x}_{ijt}$:
% \begin{description}
% \item[M1, model uczenia (PB1).] Wszystkie odpowiedzi; jedynym predyktorem jest numer próby,
%   $\beta_1 \log_2 t$. Wartość $e^{\beta_1}$ jest ilorazem szans (OR) poprawnej odpowiedzi przy
%   podwojeniu numeru próby (np.\ z 1 do 2 albo z 4 do 8).
% \item[M2, model retencji (PB2).] Wyłącznie próby powtórne ($t \geq 2$);
%   $\beta_1 \log_2 t + \beta_2 Y_{ij,t-1} + \beta_3 \log_2 \Delta_{ijt} + \beta_4 S_{it}$, gdzie
%   $Y_{ij,t-1}$ to wynik poprzedniej próby, $\Delta_{ijt}$ odstęp od niej w dniach (ograniczony od dołu
%   do 1 dnia), a $S_{it}$ staż w latach. Konstrukcją model odpowiada modelom \emph{performance factors
%   analysis} i \emph{half-life regression} stosowanym do logów uczących się \citep{pavlik2009,settles2016}.
% \item[M3, model nowych pytań.] Wyłącznie pierwsze próby ($t = 1$); jedynym predyktorem jest staż,
%   $\beta_4 S_{it}$. Przy efektach stałych pracownika współczynnik porównuje pierwsze kontakty tej samej
%   osoby z kolejnymi pytaniami, napotkanymi przy różnym stażu.
% \end{description}
% Pracowników i pytania, w których wszystkie odpowiedzi były poprawne albo wszystkie błędne, usuwano
% iteracyjnie przed estymacją, ponieważ ich efekty stałe nie są identyfikowalne. Modele szacowano metodą
% największej wiarygodności, a błędy standardowe klastrowano względem pracowników.

% Odporność oszacowania $e^{\beta_1}$ z modelu M1 sprawdzono w analizach wrażliwości: w zbalansowanej
% kohorcie, po obcięciu sekwencji do 10 prób, w warstwach długości sekwencji (kontrola obciążenia typu
% \emph{pattern-mixture}), po wykluczeniu pracowników z mniej niż 50 odpowiedziami oraz po dodaniu stażu
% jako predyktora. Efekt uśredniony populacyjnie oszacowano ponadto metodą uogólnionych równań
% estymujących (GEE) z wymienną macierzą korelacji roboczej \citep{liang1986}, a w modelu mieszanym
% efekty stałe zastąpiono skrzyżowanymi losowymi wyrazami wolnymi pracowników i pytań \citep{baayen2008}.
% Z modelu mieszanego wyznaczono współczynniki korelacji wewnątrzklasowej (ICC) w skali latentnej oraz
% brzegowy i warunkowy współczynnik $R^2$ \citep{nakagawa2013}.

% \subsection{Przejścia i powtarzane błędy}
% Przejścia między kolejnymi próbami tego samego pytania zestawiono w tabeli kontyngencji. Dla dwóch
% następujących po sobie błędów porównano odsetek ponownych wyborów tego samego dystraktora z odsetkiem
% oczekiwanym przy losowym wyborze spośród zaobserwowanych dystraktorów oraz przy wyborze proporcjonalnym
% do popularności każdego dystraktora w danym pytaniu.

% Analizy uzupełniające, czyli zmianę poprawności na poziomie pracowników oraz analizę pytań metodami
% klasycznej teorii testu, opisano w materiale uzupełniającym. Obliczenia wykonano w języku Python
% \VerPython{} z użyciem bibliotek pandas \VerPandas, NumPy \VerNumpy, SciPy \VerScipy, statsmodels
% \VerStatsmodels{} i Matplotlib \VerMatplotlib{}.
% %  (ziarno generatora liczb losowych \Seed).

\section{Wyniki}

\subsection{Pierwsza próba}
Przy pierwszym kontakcie z pytaniem (\FirstN{} odpowiedzi \NEmployees{} pracowników) poprawność
wynosiła \FirstAcc\% (95\% PU \FirstAccCI; tab.~\ref{tab:pierwsza_proba_wiek}). Poprawność
zmieniała się z wiekiem niemonotonicznie: była najniższa w przedziale \FirstAccMinGroup{} lat
(\FirstAccMin\%), a najwyższa w przedziale \FirstAccMaxGroup{} lat (\FirstAccMax\%). Przy porównaniu
odpowiedzi na te same pytania szanse poprawnej odpowiedzi rosły z wiekiem (OR \FirstAgeOR{} na 10 lat,
95\% PU \FirstAgeORCI; $p \FirstAgeP$). Bez efektów stałych pytania zależność była słabsza
(OR \FirstAgeRawOR, 95\% PU \FirstAgeRawORCI; $p \FirstAgeRawP$).

\begin{table}[H]
\centering
\caption{Poprawność w pierwszej próbie ogółem i według przedziału wieku pracownika.}
\label{tab:pierwsza_proba_wiek}
\footnotesize
\setlength{\tabcolsep}{4pt}
\tabelacsv{l*{7}{r}}{pierwsza_proba_wiek}% 7 = ogółem + przedziały wieku + brak danych
\uwagatabeli{Pierwsza próba to pierwsza odpowiedź w sekwencji. 95\% PU wyznaczono metodą bootstrap.
Zależność od wieku oszacowano regresją logistyczną opisaną w podrozdziale~\ref{sec:analiza}
(\FirstAgeN{} odpowiedzi \FirstAgeEmp{} pracowników o znanym wieku): OR \FirstAgeOR{} na 10 lat,
95\% PU \FirstAgeORCI, $p \FirstAgeP$.}
\end{table}

\subsection{Zmiany w sekwencji (PB1)}
Poprawność rosła wraz z numerem próby, z \AccAttOne\% \AccAttOneCI{} w pierwszej próbie do
\AccAttMax\% \AccAttMaxCI{} w próbie \AttMax{}. i dalszych (rys.~\ref{fig:krzywa_uczenia}). Ponieważ
z każdą próbą ubywało sekwencji, zmianę sprawdzono w kohorcie \CohortN{} sekwencji. W kohorcie
poprawność wzrosła z \CohortAccOne\% \CohortAccOneCI{} w próbie~1 do \CohortAccK\% \CohortAccKCI{}
w próbie~\CohortK{}, czyli o \SeqChange{} p.p. (95\% PU \SeqChangeCI; tab.~\ref{tab:sekwencje_wiek}).
Wzrost w tych samych sekwencjach oznacza, że poprawa nie wynikała z wcześniejszego kończenia sekwencji
o niższej poprawności. Poprawność w próbie~1 w kohorcie (\CohortAccOne\%) była zbliżona do
poprawności w pierwszej próbie we wszystkich sekwencjach (\FirstAcc\%).

Oszacowanie wzrostu było dodatnie we wszystkich przedziałach wieku i wynosiło od \SeqChangeMin{} p.p.
(\SeqChangeMinGroup{} lat) do \SeqChangeMax{} p.p. (\SeqChangeMaxGroup{} lat), przy czym w nielicznych
przedziałach skrajnych dolna granica 95\% PU była bliska zera. Nie stwierdzono zależności tempa wzrostu
od wieku (OR interakcji \SeqAgeOR{} na 10 lat wieku i podwojenie numeru próby, 95\% PU \SeqAgeORCI;
$p \SeqAgeP$).

\begin{figure}[H]
\centering
\includegraphics{\FigDir rys_krzywa_uczenia}
\caption{Poprawność według numeru próby we wszystkich sekwencjach dostępnych w danej próbie (ich liczba
maleje wraz z numerem próby). Słupki błędów: 95\% PU (bootstrap). Próby \MaxAttemptPooled{}. i dalsze połączono. Przebieg w kohorcie
przedstawiono na rys.~\ref{fig:stany_kohorty}.}
\label{fig:krzywa_uczenia}
\end{figure}

\begin{table}[H]
\centering
\caption{Poprawność w kolejnych próbach w kohorcie, ogółem i według przedziału wieku pracownika.}
\label{tab:sekwencje_wiek}
\footnotesize
\setlength{\tabcolsep}{4pt}
\tabelacsv{l*{7}{r}}{sekwencje_wiek}% 7 = ogółem + przedziały wieku + brak danych
\uwagatabeli{Kohorta: sekwencje liczące co najmniej \CohortK{} prób, próby od 1 do \CohortK{}.
p.p.: punkty procentowe; 95\% PU zmiany wyznaczono metodą bootstrap. Zależność tempa poprawy od wieku
oszacowano regresją logistyczną opisaną w podrozdziale~\ref{sec:analiza} (\SeqAgeN{} odpowiedzi
pracowników o znanym wieku).}
\end{table}

\subsection{Wynik poprzedniej próby i odstęp (PB2)}
Wzrost poprawności jest wynikiem dwóch przeciwnych procesów: korekty błędów i zapominania. Harmonogram
nie zniekształcał ich porównania, ponieważ pytanie wracało po odpowiedzi poprawnej i błędnej z podobnym
prawdopodobieństwem (\ReaskAfterCorrect\% i \ReaskAfterIncorrect\%) i po podobnym odstępie (mediana
\GapAfterCorrect{} i \GapAfterIncorrect{} dni). Po odpowiedzi poprawnej kolejna była poprawna
w \RetainPct\% \RetainCI{} przypadków, a w \ForgetPct\% wystąpiło zapominanie. Po odpowiedzi błędnej
korekta nastąpiła w \RecoverPct\% \RecoverCI{} przypadków. Odsetki te dotyczą różnych grup prób,
dlatego porównano także liczby zdarzeń: korekt błędów było \NCorrections, a przypadków zapominania
\NForgetting. Przewaga korekt wyjaśnia wzrost poprawności. Przy stałych odsetkach korekt błędów
i zapominania poprawność dążyłaby do poziomu równowagi \EquilibriumPct\%. Poprawność w kohorcie
w próbie~\CohortK{} (\CohortAccK\%) była bliska temu poziomowi, lecz w próbach \AttMax{}. i dalszych
wynosiła \AccAttMax\% i go przekraczała. Odsetek zapominania nie był bowiem stały, lecz zależał od
odstępu i od liczby poprawnych odpowiedzi z rzędu.

Odsetek zapominania rósł wraz z odstępem, z \ForgetShort\% przy odstępie krótszym niż 14 dni do
\ForgetLong\% przy odstępie co najmniej 240 dni (rys.~\ref{fig:odstep}). Odsetek korekt błędów zmieniał
się z odstępem w mniejszym stopniu (\RetShortIncorrect\% i \RetLongIncorrect\% w tych samych przedziałach).

Odsetek zapominania zależał także od liczby poprawnych odpowiedzi z rzędu. Po jednej odpowiedzi
poprawnej wynosił \ForgetRunOne\% \ForgetRunOneCI, po dwóch \ForgetRunTwo\%, po trzech
\ForgetRunThree\%, a po co najmniej \ForgetRunK{} \ForgetRunMax\% \ForgetRunMaxCI{}
(tab.~\ref{tab:seria_poprawnych}). Różnica utrzymywała się przy każdym odstępie: przy odstępie
krótszym niż 14 dni odsetek zapominania wynosił
\ForgetFreshShort\% po jednej i \ForgetStableShort\% po co najmniej trzech odpowiedziach poprawnych,
a przy odstępie co najmniej 240 dni odpowiednio \ForgetFreshLong\% i \ForgetStableLong\%. Najmniej
trwałe były korekty błędów. Po korekcie nawrót błędu wystąpił w \RelapsePct\% \RelapseCI{}
przypadków ($n = \RelapseN$), a w \RelapseSamePct\% \RelapseSameCI{} z nich pracownik wybrał ten sam
dystraktor co przed korektą. Udział ten był zbliżony do odsetka ponownych wyborów tego samego dystraktora
po dwóch kolejnych błędach (\SameDistractorPct\%; podrozdział~\ref{sec:pb3}).

\begin{figure}[H]
\centering
\includegraphics{\FigDir rys_odstep}
\caption{Poprawność według odstępu, osobno dla prób poprzedzonych odpowiedzią poprawną (koła)
i błędną (kwadraty). Dla prób poprzedzonych odpowiedzią poprawną odsetek zapominania wynosi 100\% minus
poprawność, a dla prób poprzedzonych odpowiedzią błędną poprawność jest odsetkiem korekt błędów.
Słupki błędów: 95\% PU (bootstrap).}
\label{fig:odstep}
\end{figure}

\begin{table}[H]
\centering
\caption{Odsetek zapominania według liczby poprawnych odpowiedzi udzielonych z rzędu bezpośrednio
przed daną próbą.}
\label{tab:seria_poprawnych}
\footnotesize
\setlength{\tabcolsep}{5pt}
\tabelacsv{l*{4}{r}}{seria_poprawnych}% 4 = 1, 2, 3, >=ForgetRunK
\uwagatabeli{Próby poprzedzone odpowiedzią poprawną. Odsetek zapominania: odsetek odpowiedzi błędnych.
95\% PU wyznaczono metodą bootstrap.}
\end{table}

\subsection{Utrwalone błędy (PB3)}
\label{sec:pb3}
Większość błędów korygowano: w \EverCorrectedPct\% z \FirstWrongSeqN{} sekwencji rozpoczętych
odpowiedzią błędną i liczących co najmniej dwie próby pracownik co najmniej raz odpowiedział później
poprawnie. Szansa korekty malała jednak z każdym
kolejnym błędem. Po jednym błędzie poprawność wynosiła \CorrRunOne\% \CorrRunOneCI, po dwóch błędach
z rzędu \CorrRunTwo\%, po trzech \CorrRunThree\%, a po co najmniej \RunMax{} \CorrRunMax\%
\CorrRunMaxCI{} (tab.~\ref{tab:seria_bledow}). Część tego spadku może wynikać z tego, że kilka błędów
z rzędu zdarzało się głównie w trudniejszych pytaniach.

Powtórzone błędy nie były zgodne z wyborem losowym. Po dwóch kolejnych błędach pracownik w \SameDistractorPct\%
\SameDistractorCI{} przypadków ($n = \SameDistractorN$) wybrał ten sam dystraktor, podczas gdy przy
wyborze losowym należałoby oczekiwać \SameDistractorExpected\%, a przy wyborze zgodnym z popularnością
dystraktorów \SameDistractorExpectedPop\%. Odsetek ten nie zależał od odstępu (\SameDistrGapShort\%
przy odstępie krótszym niż 30 dni, \SameDistrGapMid\% przy 30--89 dniach i \SameDistrGapLong\% przy
co najmniej 90 dniach). Powtórzone błędy skupiały się w niewielu pytaniach: na \PersistTopItems{}
ze \PersistItems{} pytań przypadało \PersistTopItemsShare\% z nich.

\begin{table}[H]
\centering
\caption{Poprawność według liczby błędów popełnionych z rzędu bezpośrednio przed daną próbą.}
\label{tab:seria_bledow}
\footnotesize
\setlength{\tabcolsep}{5pt}
\tabelacsv{l*{5}{r}}{seria_bledow}% 5 = 0, 1, 2, 3, >=RunMax
\uwagatabeli{Próby od drugiej wzwyż. Kolumna 0: próby poprzedzone odpowiedzią poprawną. 95\% PU
wyznaczono metodą bootstrap.}
\end{table}

\subsection{Podsumowanie wyników}
Rys.~\ref{fig:stany_kohorty} zestawia wyniki w kohorcie. Poprawność rośnie z próby na próbę (PB1),
korekty błędów przeważają nad zapominaniem (PB2), a w każdej próbie od drugiej wzwyż część odpowiedzi
błędnych stanowią powtórzone błędy (PB3).

\begin{figure}[H]
\centering
\includegraphics{\FigDir rys_stany_kohorty}
\caption{Poprawność i stan odpowiedzi w kohorcie (\CohortN{} sekwencji, próby od 1 do \CohortK{}).
Słupki: udział odpowiedzi według wyniku bieżącej i poprzedniej próby (podrozdział~\ref{sec:analiza});
w próbie~1 tylko odpowiedzi poprawne i błędne. Linia: poprawność z 95\% PU (bootstrap),
wartości nad słupkami.}
\label{fig:stany_kohorty}
\end{figure}

% \section{Wyniki - stare do usuniecia}

% \subsection{Próba i harmonogram powtórek}
% Zbiór danych obejmował \NResponses{} odpowiedzi \NEmployees{} pracowników na \NItems{} pytań
% w \NModules{} modułach, zebranych od \DateStart{} do \DateEnd{} (tab.~\ref{tab:zbior_danych}).
% Mediana liczby odpowiedzi na pracownika wynosiła \RespPerEmpMedian{} (IQR
% \RespPerEmpQOne--\RespPerEmpQThree), udzielonych na \ItemsPerEmpMedian{} różnych pytań w ciągu
% \ActiveDaysMedian{} dni aktywności, a mediana czasu uczestnictwa w programie wynosiła
% \SpanMonthsMedian{} miesięcy. Okresy bez odpowiedzi odpowiadają zawieszeniom programu na czas
% przeglądu pytań (rys.~\ref{fig:aktywnosc}). Wśród \DemoMatched{} pracowników o znanych danych
% demograficznych kobiety stanowiły \DemoWomenPct\%, a mediana wieku wynosiła \AgeMedian{} lat
% (tab.~\ref{tab:demografia}). Ogólny odsetek poprawnych odpowiedzi wynosił \AccOverall\%, a w pierwszych
% próbach \AccFirst\%.

% Harmonogram powtórek nie zależał od wyniku poprzedniej próby: prawdopodobieństwo kolejnej próby
% wynosiło \ReaskAfterCorrect\% po odpowiedzi poprawnej i \ReaskAfterIncorrect\% po błędnej, natomiast
% mediana odstępu do kolejnej próby wynosiła odpowiednio \GapAfterCorrect{} i \GapAfterIncorrect{} dni.
% Kontynuacja sekwencji nie była więc selekcjonowana według wyników. % SPRAWDŹ

% \subsection{Uczenie się w kolejnych próbach (PB1)}
% Odsetek poprawnych odpowiedzi rósł wraz z numerem próby: z \AccAttOne\% \AccAttOneCI{} w pierwszej
% próbie do \AccAttMax\% \AccAttMaxCI{} w próbie \AttMax{}. i kolejnych (rys.~\ref{fig:krzywa_uczenia},
% tab.~\ref{tab:krzywa_uczenia}). Wzrost był podobny w zbalansowanej kohorcie \CohortN{} sekwencji
% liczących co najmniej \CohortK{} prób (\CohortAccOne\% \CohortAccOneCI{} w próbie~1 wobec \CohortAccK\%
% \CohortAccKCI{} w próbie~\CohortK), nie wynika więc z selektywnej kontynuacji przez pracowników
% osiągających lepsze wyniki. Odsetek błędów był funkcją potęgową numeru próby (wykładnik \PowerLawB,
% $R^2 = \PowerLawRsq$), zgodnie z potęgowym prawem praktyki \citep{newell1981}.

% W modelu uczenia M1 każde podwojenie numeru próby wiązało się z \ORMOneAttempt-krotnie wyższymi
% szansami poprawnej odpowiedzi (95\% PU \ORMOneAttemptCI; tab.~\ref{tab:regresja}). Efekt ten wyjaśniał
% jednak niewielką część zmienności poprawności (brzegowy $R^2 = \RsqMarginal$ wobec warunkowego
% $R^2 = \RsqConditional$), a pytania odpowiadały za większą jej część niż pracownicy (ICC
% \ICCItem{} wobec \ICCEmployee). Na poziomie pracowników wynik w trzeciej próbie był lepszy niż
% w pierwszej u \EmpImproved{} z \EmpN{} osób, a gorszy u \EmpWorsened{} (średnia zmiana
% \EmpMeanChange{} p.p., 95\% PU \EmpMeanChangeCI; tab.~\ref{tab:zmiana_pracownikow}).

% \begin{table*}[htbp]
% \centering
% \caption{Modele regresji logistycznej poprawności odpowiedzi: ilorazy szans [95\% PU].}
% \label{tab:regresja}
% \small
% \setlength{\tabcolsep}{3.5pt}
% \tabelacsv{llll}{regresja}
% \uwagatabeli{Definicje modeli M1--M3 podano w podrozdziale ,,Modele statystyczne''. Wszystkie modele
% zawierają efekty stałe pracownika i pytania; błędy standardowe klastrowane względem pracowników.
% Pracowników i pytania bez zmienności poprawności usunięto przed estymacją (M1: \MOneDropped{}
% odpowiedzi). OR dla numeru próby i odstępu dotyczą podwojenia wartości, dla stażu jednego roku.}
% \end{table*}

% \subsection{Wynik poprzedniej próby, odstęp i zapominanie (PB2)}
% W modelu retencji M2, po uwzględnieniu numeru próby, poprawna poprzednia odpowiedź wiązała się
% z wyższymi szansami poprawnej odpowiedzi (OR \ORMTwoPrev, 95\% PU \ORMTwoPrevCI), natomiast dłuższy
% odstęp od poprzedniej próby wiązał się z niższymi szansami (OR \ORMTwoGap{} na podwojenie odstępu,
% 95\% PU \ORMTwoGapCI; tab.~\ref{tab:regresja}). Wśród odpowiedzi poprzednio poprawnych odsetek
% poprawnych zmniejszał się z \RetShortCorrect\%, gdy pytanie powtórzono w ciągu 14 dni, do
% \RetLongCorrect\% po 240 dniach lub dłużej; odpowiedzi poprzednio błędne były korygowane odpowiednio
% w \RetShortIncorrect\% i \RetLongIncorrect\% przypadków (rys.~\ref{fig:odstep},
% tab.~\ref{tab:retencja_wg_odstepu}). Wzorzec ten jest zgodny z zapominaniem między powtórkami
% \citep{cepeda2006,cepeda2008}. % SPRAWDŹ

% \subsection{Przejścia i utrwalone błędy (PB3)}
% Spośród odpowiedzi poprzednio poprawnych \RetainPct\% \RetainCI{} było poprawnych również w kolejnej
% próbie, a \ForgetPct\% zapomniano. Spośród odpowiedzi poprzednio błędnych \RecoverPct\% \RecoverCI{}
% skorygowano, a \PersistPct\% pozostało błędnych, mimo że po każdej odpowiedzi wyświetlano poprawną
% odpowiedź (tab.~\ref{tab:przejscia}). Gdy błąd się powtarzał (\SameDistractorN{} przypadków),
% pracownicy wybierali ten sam dystraktor w \SameDistractorPct\% \SameDistractorCI{} przypadków, wobec
% \SameDistractorExpected\% oczekiwanych przy wyborze losowym i \SameDistractorExpectedPop\% oczekiwanych
% przy uwzględnieniu popularności poszczególnych dystraktorów. Wskazuje to na utrwalone błędne
% przekonania, a nie na zgadywanie. % SPRAWDŹ

% \subsection{Odporność wyników i nowe pytania}
% Oszacowanie efektu numeru próby pozostawało zbliżone we wszystkich specyfikacjach (OR od \ORRange;
% rys.~\ref{fig:wrazliwosc}, tab.~\ref{tab:wrazliwosc}) i wynosiło \ORSensCohort{} w zbalansowanej
% kohorcie, \ORSensLenThree, \ORSensLenFive{} i \ORSensLenEight{} w sekwencjach liczących 3--4, 5--7
% i co najmniej 8 prób, \ORSensGee{} w uśrednionym populacyjnie modelu GEE oraz \ORSensMixed{} w modelu
% mieszanym. Oszacowanie było wyższe po dodaniu stażu jako predyktora (\ORSensTenure), ponieważ staż
% jest skorelowany z numerem próby.

% W modelu nowych pytań M3 odsetek poprawnych odpowiedzi przy pierwszym kontakcie z pytaniem nie rósł
% wraz ze stażem (OR \ORMThreeTenure{} na rok, 95\% PU \ORMThreeTenureCI; tab.~\ref{tab:regresja}).
% Przedział ufności obejmuje zarówno brak efektu, jak i niewielki spadek, dlatego wynik oznacza brak
% dowodu na uogólnienie poprawy na nowe pytania, a nie dowód jego braku. % SPRAWDŹ

% \begin{figure}[tbp]
% \centering
% \includegraphics{\FigDir rys_wrazliwosc}
% \caption{Iloraz szans (95\% PU) poprawnej odpowiedzi na podwojenie numeru próby w różnych
% specyfikacjach modelu M1. Romb: model główny; linia pionowa: brak efektu.}
% \label{fig:wrazliwosc}
% \end{figure}

% =============================================================================
% Dyskusja i wnioski. Plik edytowany ręcznie -- skrypt analiza_odpowiedzi_naukowa.py go NIE nadpisuje.
% Liczby: definicje.tex; literatura: references.bib.
% =============================================================================
\section{Dyskusja}

W programie mikrouczenia w banku poprawność odpowiedzi rosła z każdym powrotem pytania, także w kohorcie, w której
porównywano te same sekwencje. Wzrost wynikał z przewagi korekt błędów nad zapominaniem. Część błędów
utrzymywała się jednak mimo informacji zwrotnej i polegała na ponownym wyborze tego samego dystraktora.
Kierunek zmian jest zgodny z badaniami edukacji rozłożonej w czasie w kształceniu medycznym.
U lekarzy specjalizujących się w urologii przewaga wiedzy uzyskana dzięki cyklicznie powracającym
pytaniom utrzymywała się jeszcze dwa lata po zakończeniu programu \citep{kerfoot2009}, a przegląd
systematyczny badań nad kształceniem ustawicznym klinicystów wskazał na korzystny wpływ tej metody na
przyswajanie i utrwalanie wiedzy \citep{phillips2019}.

Warunki tamtych badań różniły się jednak od warunków niniejszej analizy. Uczestnikami byli pracownicy
ochrony zdrowia, a efekt oceniano zazwyczaj testem przeprowadzanym poza programem, w porównaniu
z grupą kontrolną. W badanym programie uczestniczyli pracownicy banku, a wzrost poprawności
obserwowano wyłącznie w odpowiedziach udzielanych w samym programie. Zgodność dotyczy zatem kierunku
zmian, nie zaś wielkości efektu ani jego przyczyny. Uzyskane wyniki pokazują natomiast, z czego
składał się wzrost poprawności, czego nie ujawnia wynik testu końcowego. Był on różnicą dwóch
przeciwnych procesów, korekt błędów i zapominania (\NCorrections{} wobec \NForgetting{} zdarzeń),
a ten sam łączny wzrost mógłby powstać także przy częstszym zapominaniu, jeśli towarzyszyłaby mu
odpowiednio większa liczba korekt. Miary zagregowane nie pozwalają też wykryć błędów, które
utrzymują się mimo informacji zwrotnej.

Zapominanie wiązało się z trzema czynnikami. Pierwszym był upływ czasu: zapominanie było częstsze przy
dłuższym odstępie, zgodnie z wynikami badań nad odstępem między powtórkami \citep{cepeda2006}. Drugim
była krótka historia poprawnych odpowiedzi. Wiedza potwierdzona jedną poprawną odpowiedzią była
zapominana od dwóch do czterech razy częściej niż wiedza potwierdzona kilka razy z rzędu, przy każdym
odstępie.
Podobnie w badaniach nad wielokrotnym powtarzaniem materiału każde kolejne udane przypomnienie
zwiększało trwałość wiedzy \citep{rawson2013}. Trzecim czynnikiem był nawrót błędu. Korekta błędu
okazała się najmniej trwała, a jeśli się nie utrzymała, pracownik zwykle wracał do dystraktora
wybranego przed korektą. Powrót pierwotnych błędów po informacji zwrotnej opisano również w badaniach
laboratoryjnych \citep{butler2011}. Poprawna odpowiedź mogła więc nie zastępować błędnego przekonania,
lecz współistnieć z nim.

Powtórzone błędy miały inny charakter niż zapominanie. Pracownicy wybierali ten sam dystraktor
częściej, niż wynikałoby z przypadku, niezależnie od odstępu, a każdy kolejny błąd zmniejszał szansę
korekty. Taki wzorzec jest zgodny z utrwalonym błędnym przekonaniem, a nie z brakiem wiedzy, który
prowadziłby do zgadywania. Wybrany dystraktor może utrwalić się w pamięci, a informacja zwrotna
ogranicza ten efekt \citep{butler2008}. W badanym programie informacja zwrotna podawała poprawną
odpowiedź i wyjaśnienie, ale nie odnosiła się do wybranego dystraktora, choć wyjaśnienie, dlaczego
wybrana odpowiedź jest błędna, działa zwykle silniej \citep{shute2008}. Skupienie powtórzonych błędów
w \PersistTopItems{} pytaniach sugeruje ponadto, że część z nich wynikała z konstrukcji pytania, np.\
z niejednoznacznego dystraktora lub treści niezgodnej z praktyką banku. Takie pytania można wskazać
standardową analizą dystraktorów \citep{gierl2017}.

W szkoleniach jednorazowych, np.\ z bezpieczeństwa informacji, efekty zanikały po kilku miesiącach
\citep{reinheimer2020}. Uzyskane wyniki wskazują, że wiedza wymaga dalszych powtórek także w programie
z cyklicznie powracającymi pytaniami, a sposób planowania tych powtórek można poprawić. Ponieważ odstęp
nie zależał od wyniku, pytania opanowane wracały równie często jak pytania, w których pracownik dopiero
skorygował błąd. Harmonogram adaptacyjny mógłby wydłużać odstęp po kilku poprawnych odpowiedziach
z rzędu i skracać go po korekcie błędu, gdy ryzyko nawrotu jest największe. Personalizacja powtórek
poprawiała długotrwałe zapamiętywanie w porównaniu z harmonogramem stałym \citep{lindsey2014}, a modele
tempa zapominania oparte na historii odpowiedzi stosuje się już w aplikacjach do nauki języków
\citep{settles2016}. Drugą zmianą może być informacja zwrotna odnosząca się do wybranego dystraktora,
wprowadzona najpierw w pytaniach z największą liczbą powtórzonych błędów.

\subsection{Ograniczenia i dalsze badania}
Badanie było obserwacyjne i nie miało grupy kontrolnej, dlatego wzrostu poprawności nie można przypisać
wyłącznie programowi, ponieważ pracownicy zdobywali wiedzę także w codziennej pracy. Ponieważ pytania
się powtarzały, dane nie pozwalają odróżnić utrwalenia wiedzy od zapamiętania odpowiedzi na konkretne
pytanie. Nie rejestrowano czasu namysłu podczas udzielania odpowiedzi ani pewności pracownika, dlatego zapominania
nie można oddzielić od nieuwagi ani od trafnego zgadnięcia we wcześniejszej próbie, które modele ucznia
ujmują jako poślizg i zgadywanie \citep{corbett1994}. Na takie źródło części zapominania wskazuje jego
odsetek wynoszący \ForgetShort\% już przy odstępie krótszym niż 14 dni. Najdłuższe odstępy mogły
obejmować okresy zawieszenia programu, dlatego zmiany treści pytań mogły zawyżać odsetek zapominania
przy tych odstępach. Porównania miały charakter opisowy i nie korygowano ich ze względu na
wielokrotność testów. Przedziały ufności uwzględniały zależność odpowiedzi tego samego pracownika,
lecz nie tego samego pytania. Zależności od liczby poprawnych odpowiedzi lub błędów z rzędu mogą więc
częściowo odzwierciedlać różnice trudności między pytaniami.

Dalsze badania mogłyby porównać w schemacie randomizowanym harmonogram stały z adaptacyjnym oraz
informację zwrotną standardową z odnoszącą się do wybranego dystraktora, z użyciem równoważnych wersji
pytań, które oddzieliłyby wiedzę od zapamiętania odpowiedzi. Rejestrowanie pewności odpowiedzi
pozwoliłoby sprawdzić, czy utrwalone błędy dotyczą przekonań o wysokiej pewności \citep{butler2011}.
Wyniki mikrouczenia warto też powiązać z zastosowaniem wiedzy w pracy, np.\ z liczbą błędów
proceduralnych \citep{blume2010}.

\section{Wnioski}
W programie z cyklicznie powracającymi pytaniami poprawność odpowiedzi pracowników banku rosła z próby
na próbę. Wiedza wymagała jednak dalszych powtórek, ponieważ była zapominana z upływem czasu, zwłaszcza
po pierwszej poprawnej odpowiedzi i po korekcie błędu. Część błędów utrzymywała się mimo informacji
zwrotnej i polegała na ponownym wyborze tego samego dystraktora. Harmonogram dopasowany do wyników
pracownika i informacja zwrotna odnosząca się do wybranego dystraktora mogą zwiększyć skuteczność
takich programów, ale ich wpływ wymaga potwierdzenia w badaniach z grupą kontrolną.


\section*{Oświadczenia}

% SPRAWDŹ: sposób kodowania danych pracowników oraz status zgody
% komisji etycznej lub zwolnienia z obowiązku jej uzyskania.
\noindent\textbf{Etyka.}
W badaniu wykorzystano w pełni anonimizowane dane przekazane przez bank za jego zgodą
pisemną w dniu 3 sierpnia 2026. Sposób oznaczenia pracowników i zasady ochrony danych opisano
w podrozdziale~\ref{sec:dane}.\\
\textbf{Dostępność danych.}
Dane wykorzystane w badaniu stanowią własność banku i nie są
udostępniane publicznie.\\
\textbf{Finansowanie.}
Badanie nie otrzymało finansowania.\\
\textbf{Konflikt interesów.}
M. Szerszeń jest pracownikiem banku, w którym prowadzono program badań. Pozostali autorzy deklarują
brak konfliktu interesów.\\
\textbf{Wykorzystanie narzędzi AI.}
Wykorzystanie narzędzi AI. W trakcie przygotowywania artykułu wykorzystano narzędzie Antropic Claude Sonnet 5.5, do (1) redakcji językowej opisu wyników analizy przeprowadzonej przez autorów oraz (2) tłumaczenia tekstu z języka polskiego na angielski. Analiza danych, interpretacja wyników i wnioski zostały wykonane przez autora. Autorzy zweryfikowli zredagowali całość tekstu i ponoszą pełną odpowiedzialność za jego treść.

\bibliography{references}
\end{document}
